\documentclass[10pt, a3paper, landscape, twocolumn]{article}

%----------------------------------------------------------------------------------------
%	宏包引入
%----------------------------------------------------------------------------------------
\usepackage[UTF8]{ctex}       % 中文支持
\usepackage{geometry}         % 页面设置
\usepackage{amsmath, amssymb} % 数学公式
\usepackage{graphicx}         % 图片
\usepackage{fancyhdr}         % 页眉页脚
\usepackage{tabularx}         % 自动宽度表格
\usepackage{tikz}             % 绘图（用于密封线）
\usepackage{eso-pic}          % 背景图
\usepackage{enumitem}         % 列表定制
\usepackage{lastpage}         % 用于计算总页数
\usepackage{calc}             % 用于页码计算

%----------------------------------------------------------------------------------------
%	页面设置 (A3 横向 双栏)
%----------------------------------------------------------------------------------------
\geometry{
    left=40mm,      % 左侧留出装订线空间
    right=15mm,     
    top=20mm,       
    bottom=25mm,    % 下边距稍微加大，给页码留空间
    columnsep=20mm  % 中缝宽度
}

%----------------------------------------------------------------------------------------
%	背景绘制：装订线 + 中缝线
%----------------------------------------------------------------------------------------
\AddToShipoutPictureBG{%
    \begin{tikzpicture}[remember picture, overlay]
        % 1. 左侧密封线
        \draw[dashed, thick] ([xshift=25mm, yshift=-10mm]current page.north west) -- ([xshift=25mm, yshift=10mm]current page.south west);
        \node[rotate=90, anchor=center, align=center, text width=20cm] at ([xshift=15mm]current page.west) {
            \small \kaishu 
            考生姓名\underline{\hspace{3cm}} \hspace{1cm} 
            学号\underline{\hspace{3cm}} \hspace{1cm} 
            \dots\dots\dots\dots 装 \dots\dots\dots\dots 订 \dots\dots\dots\dots 线 \dots\dots\dots\dots
        };
        % 2. 中间分隔线
        \draw[dashed, help lines] (current page.north) -- (current page.south);
    \end{tikzpicture}
}

%----------------------------------------------------------------------------------------
%	页眉页脚设置 (核心修改部分)
%----------------------------------------------------------------------------------------
\pagestyle{fancy}
\fancyhf{} % 清空默认设置

% --- 页眉 ---
\fancyhead[C]{\Large\bfseries 弘仕专升本教育 2025-2026 学年高等数学课程考核试卷}

% --- 页脚 (实现左右分栏独立页码) ---
\fancyfoot[C]{%
    % 左栏页码容器 (占据页面左半部分)
    \begin{minipage}[b]{0.45\textwidth}
        \centering
        % 计算公式：物理页码 * 2 - 1
        第 \the\numexpr\value{page}*2-1\relax\ 页 \ / \ 共 4 页
    \end{minipage}%
    \hfill % 把左右两个容器撑开
    % Right Column Page Number Container
    \begin{minipage}[b]{0.45\textwidth}
        \centering
        % 计算公式：物理页码 * 2
        第 \the\numexpr\value{page}*2\relax\ 页 \ / \ 共 4 页
    \end{minipage}%
}

\renewcommand{\headrulewidth}{0pt}

%----------------------------------------------------------------------------------------
%	列表与表格配置
%----------------------------------------------------------------------------------------
\setlist[enumerate,1]{label=\textbf{\arabic*.}, itemsep=0.5em, leftmargin=*}

\begin{document}

%----------------------------------------------------------------------------------------
%	试卷抬头
%----------------------------------------------------------------------------------------

\noindent
\textbf{试卷}：A \quad \textbf{方式}：闭卷 \quad \textbf{时间}：120 分钟 \quad \textbf{学院}：数学与统计学院

\vspace{0.1cm}

\begin{center}
    \large \textbf{XX 专业 20XX 级《高等数学》课程考核试卷}
\end{center}

\vspace{0.1cm}

% --- 考生须知框 (宽度 0.9\linewidth 且居中) ---
\begin{center}
    \fbox{%
        \parbox{\dimexpr 0.9\linewidth - 2\fboxsep - 2\fboxrule \relax}{
            \centering\textbf{考生须知}
            \begin{enumerate}[label=\arabic*., nosep, leftmargin=2em, labelsep=0.5em]
                \small
                \item 考生需携带有效身份证件参加考试，无证件者须到所在学院开具证明。
                \item 迟到 30 分钟以上者不准进入考场；考核开始 30 分钟后，考生方可交卷离场。
                \item 禁止携带手机等与考试无关的违禁物品进入考场。
                \item 严格遵守考核纪律，违者按学校相关规定处理。
            \end{enumerate}
        }%
    }
\end{center}

\vspace{0.2cm}

% --- 评分栏 (宽度 0.9\linewidth，与上表对齐) ---
\begin{center}
    \begin{tabularx}{0.9\linewidth}{| *{6}{>{\centering\arraybackslash}X|} }
        \hline
        题号 & 一 & 二 & 三 & 总分 & 复核 \\
        \hline
        分值 & 60 & 20 & 70 & 150 & \\
        \hline
        得分 & & & & & \\
        \hline
    \end{tabularx}
\end{center}

\vspace{0.2cm}

%----------------------------------------------------------------------------------------
%	一、选择题
%----------------------------------------------------------------------------------------

\section*{一、选择题：本大题共 12 小题，每小题 5 分，共 60 分.}

\begin{enumerate}
    % Q1
    \item 函数 $f(x) = \arcsin \frac{x+1}{2} + \frac{1}{\sqrt{x^2-1}}$ 的定义域为 ( \quad ) \\
    \begin{tabularx}{\linewidth}{@{} *{4}{X} @{}}
        A. $[-3,-1]$ & B. $(-3,-1]$ & C. $[-3,-1)$ & D. $(-3,-1)$
    \end{tabularx}

    % Q2
    \item 当 $x \to 0$ 时，与 $x$ 等价的无穷小量是 ( \quad ) \\
    \begin{tabularx}{\linewidth}{@{} *{2}{X} @{}}
        A. $\sqrt[3]{x} + x$ & B. $1 - e^{x^2}$ \\
        C. $\sin x$        & D. $2(1-\cos x)$
    \end{tabularx}

    % Q3
    \item 设 $f'(x_0)$ 存在，则下列 4 个式子中等于 $f'(x_0)$ 的是 ( \quad ) \\
    \begin{tabularx}{\linewidth}{@{} X @{}}
        A. $\lim\limits_{n \to \infty} n \left[ f\left(x_0 - \frac{1}{n}\right) - f(x_0) \right]$ \\
        B. $\lim\limits_{x \to x_0} \frac{f(x) - f(x_0)}{x - x_0}$ \\
        C. $\lim\limits_{\Delta x \to 0} \frac{f(x_0 + \Delta x) - f(x_0 - \Delta x)}{\Delta x}$ \\
        D. $\lim\limits_{\Delta x \to 0} \frac{f(x_0 + 2\Delta x) - f(x_0 - 3\Delta x)}{\Delta x}$
    \end{tabularx}

    % Q4
    \item 设函数 
    $f(x) = \begin{cases} 
    \frac{2019}{x} \sin x + x \sin \frac{2020}{x}, & x < 0 \\
    a - 1, & x = 0 \\
    (1 - x)^{\frac{b}{x}}, & x > 0
    \end{cases}$ 
    在 $(-\infty, +\infty)$ 内连续，则 $a = $ ( \quad ) \\
    \begin{tabularx}{\linewidth}{@{} *{1}{X} @{}}
        A. $2021 + \ln 2020$ \\
        B. $2021 - \ln 2020$ \\
        C. $2020 + \ln 2019$ \\
        D. $2020 - \ln 2019$
    \end{tabularx}

    % Q5
    \item 下列结论错误的是 ( \quad ) \\
    \begin{enumerate}[label=\Alph{enumii}., leftmargin=2em]
        \item 如果函数 $f(x)$ 在 $x=x_0$ 处连续，则 $f(x)$ 在 $x=x_0$ 处可微
        \item 如果函数 $f(x)$ 在 $x=x_0$ 处不连续，则 $f(x)$ 在 $x=x_0$ 处不可微
        \item 如果函数 $f(x)$ 在 $x=x_0$ 处可微，则 $f(x)$ 在 $x=x_0$ 处连续
        \item 如果函数 $f(x)$ 在 $x=x_0$ 处不可微，则 $f(x)$ 在 $x=x_0$ 处也可能连续
    \end{enumerate}

    % Q6
    \item 设函数 $f(x) = x + \frac{4}{x}$，其单调递增的区间 ( \quad ) \\
    \begin{tabularx}{\linewidth}{@{} *{2}{X} @{}}
        A. $(-\infty,-2),(2,+\infty)$ & B. $(-2,0) \cup (0,2)$ \\
        C. $(-\infty,0),(0,+\infty)$  & D. $(-2,0),(0,2)$
    \end{tabularx}

    % Q7
    \item 设函数 $f(x) = x^3 - 6x^2 + 9x + 2$ 的拐点是 ( \quad ) \\
    \begin{tabularx}{\linewidth}{@{} *{4}{X} @{}}
        A. $(1,6)$ & B. $(2,3)$ & C. $(2,4)$ & D. $(3,2)$
    \end{tabularx}

    % Q8
    \item 设 $f(x) = e^{-2x}$，则 $\int \frac{f'(\ln x)}{x} dx =$ ( \quad ) \\
    \begin{tabularx}{\linewidth}{@{} *{2}{X} @{}}
        A. $2x + C$ & B. $-2x + C$ \\
        C. $\frac{1}{x^2} + C$ & D. $-\frac{1}{x^2} + C$
    \end{tabularx}

    % Q9
    \item $\int_{-1}^{1} (|x| + \sin x \cos^2 x) dx =$ ( \quad ) \\
    \begin{tabularx}{\linewidth}{@{} *{4}{X} @{}}
        A. $1$ & B. $2$ & C. $-1$ & D. $-2$
    \end{tabularx}

    % Q10
    \item 过点 $P_0(4,3,1)$ 且与平面 $3x+2y+5z-1=0$ 垂直的直线方程为 ( \quad ) \\
    \begin{tabularx}{\linewidth}{@{} X @{}}
        A. $\frac{x-4}{3} = \frac{y-3}{2} = \frac{z-1}{5}$ \\
        B. $3x+2y+5z-23=0$ \\
        C. $\frac{x-4}{-3} = \frac{y-3}{17} = \frac{z-1}{-5}$ \\
        D. $3x-17y+5z+34=0$
    \end{tabularx}

    % Q11
    \item 下列所给级数中：
    (1) $\sum\limits_{n=1}^\infty \frac{n}{n+1}$, \quad
    (2) $\sum\limits_{n=1}^\infty \frac{(-1)^n}{n}$, \quad
    (3) $\sum\limits_{n=1}^\infty \ln^n 2$, \quad
    (4) $\sum\limits_{n=1}^\infty \frac{n+1}{2^n}$ \\
    收敛级数的个数为 ( \quad ) \\
    \begin{tabularx}{\linewidth}{@{} *{4}{X} @{}}
        A. 1 & B. 2 & C. 3 & D. 4
    \end{tabularx}

    % Q12
    \item 微分方程 $y' - y\sin x = e^{-\cos x}$ 的通解为 ( \quad ) \\
    \begin{tabularx}{\linewidth}{@{} *{1}{X} @{}}
        A. $y = e^{\cos x}(x+C)$ \\
        B. $y = e^{-\cos x}(x+C)$ \\
        C. $y = e^{\cos x}(-x+C)$ \\
        D. $y = e^{-\cos x}(-x+C)$
    \end{tabularx}
\end{enumerate}

%----------------------------------------------------------------------------------------
%	二、填空题
%----------------------------------------------------------------------------------------
\section*{二、填空题: 每小题 5 分, 共 20 分.}

\begin{enumerate}[start=13]
    \item 设 $y = f(x)$ 满足 $\ln y + \frac{x}{y} = 1$，则 $dy = \underline{\hspace{2cm}}$.
    \item 幂级数 $\sum\limits_{n=1}^\infty \frac{(-1)^n}{\sqrt{n}}(x-5)^n$ 的收敛域为 $\underline{\hspace{2cm}}$.
    \item 椭圆抛物面 $z = 2x^2 + 3y^2 - 5$ 在点 $(2,-1,6)$ 处的法线方程为 $\underline{\hspace{2cm}}$.
    \item 定积分 $\int_{-1}^{1} (x^6 - x^2 + 2x^3\sqrt{1-x^2} + 1) dx = \underline{\hspace{2cm}}$.
\end{enumerate}

% 强制换栏（从这里开始是 A3 纸的右半部分，即逻辑上的“第 2 页”）
\newpage 

%----------------------------------------------------------------------------------------
%	三、计算题
%----------------------------------------------------------------------------------------
\section*{三、计算题：共 70 分. 须写出计算过程.}

\begin{enumerate}[start=17]
    % Q17
    \item (5分) 求极限 $\lim\limits_{x \to 0} \frac{e^x - e^{-x} - 2x}{x - \sin x}$.
    \vspace{4cm}

    % Q18
    \item (10分) 求不定积分 $\int x \tan^2 x \, dx$.
    \vspace{4cm}

    % Q19
    \item (15分) 设 $z = x^2 f\left(\frac{y}{x}, x+y\right)$，其中 $f$ 具有连续的二阶偏导数，求 $\frac{\partial z}{\partial x}, \frac{\partial^2 z}{\partial x \partial y}$.
    \vspace{6cm}

    % Q20
    \item (10分) 求微分方程 $y' + 2xy = 4x$ 满足 $y|_{x=0} = 1$ 时的特解.
    \vspace{6cm}

    % Q21
    \item (15分) 计算二重积分 $\iint_{D} (a - \sqrt{x^2+y^2}) \, d\sigma$，其中区域 $D$ 由 $x^2 + y^2 = a^2 \, (a > 0)$ 所围成.
    \vspace{6cm}

    % Q22
    \item (15分) 要制作一个体积为 $576 \, \text{cm}^3$ 的长方体带盖的盒子，其底面长宽之比为 $2:1$，问长、宽、高各取何值时，才能使盒子的表面积最小？
    \vspace{6cm}
\end{enumerate}

\end{document}